\section{Thermal orbit weights: the current result}
\label{sec:scope}

The current investigation seeks an independently defined physical system
whose response supplies the arithmetic side of the shifted-zeta
differential equation. Its most useful new result concerns the changing
orbit weights. The Bost--Connes system has native projections that detect
primitiveness at specified primes. Their positive thermal expectations,
with one common inverse temperature, give precisely the product entering
every required arithmetic coefficient. No separate choice is made for
each prime or each repeated orbit.

More precisely, let $\ell$ always denote a prime, let $\varphi_\beta$ be
a Kubo--Martin--Schwinger (KMS) equilibrium state at inverse temperature
$\beta$, and let $P_n$ select states primitive at every prime dividing
$n$. The dictionary is
\begin{equation}
 \boxed{\quad \beta=2\omega,\qquad
 c_n(\omega)=n^{(\beta-1)/2}\varphi_\beta(P_n),\qquad
 \varphi_\beta(P_n)=\prod_{\ell\mid n}(1-\ell^{-\beta}).\quad}
 \label{eq:main_dictionary}
\end{equation}
The primitive projector is independent of temperature. The displayed
centering factor is additional and explicit; its physical origin is not
claimed to follow merely from the existence of a thermal state.

The established Bost--Connes theory supplies positive KMS states in the
range $0<\beta\le1$, although the usual global Gibbs trace does not
converge there \cite{BC,Neshveyev}. This is the range needed to move from
the modular half shift $\omega=1/2$ toward zero. Thus the candidate is
more informative than a formal identification of a zeta partition
function. A finite primitive observable has a meaningful positive
expectation at the relevant temperatures.

The contribution of this investigation is the exact comparison with
the desired orbit coefficients and the analysis of specific readout and
norm requirements. The underlying equilibrium system, KMS existence
theorem and primitive-count identities are established mathematics.
There is no claim of literature novelty for them.

Three levels must remain separate. First, the primitive probabilities
in \eqref{eq:main_dictionary} are native thermal observables. Second,
their centered values reproduce the Euler coefficients, including the
entire prime part of the shift generator. Third, a causal physical
response with these amplitudes, the archimedean completion and the
ordinary radiation norm has not been constructed. Positivity at the
first level does not prove the third.

The completed interface tests sharpen this distinction. A controlled
unitary delay model realizes actual logarithmic flight times, an ordinary
energy balance and an identity limit, but has the wrong arithmetic
amplitudes. Moreover, the unchanged archimedean factor times a finite
Euler product cannot be globally passive. A subsequent energy-exchanging
thermal load retains the full modular core and is passive as a linear
boundary system, yet fails a structural front-exponent test. These
results justify parking this regular boundary-load branch, not rejecting
the equilibrium coefficient identity or every possible global coupling.

Sections~\ref{sec:target}--\ref{sec:thermal} establish the dictionary,
and Section~\ref{sec:norms} gives the local norms and assembly limitations.
Section~\ref{sec:interface} presents the completed finite-place test and
completion obstruction; Section~\ref{sec:thermal_boundary} gives the
energy-exchanging response and finite-spectral-mass theorem.
Section~\ref{sec:previous} summarizes earlier tests, and
Section~\ref{sec:outlook} states the remaining priorities. The principal
research sources are the orbit-weight note \cite{OrbitNote} and the two
subsequent test notes \cite{InterfaceNote,BoundaryNote}. Their numerical
records and same-assistant audits are distinguished from the analytical
arguments in Appendix~\ref{app:validation}.
